\section{Wstęp}

\paragraph{}
Bezzałogowe statki powietrzne (UAV – \textit{Unmanned Aerial Vehicle}) stanowią najważniejszą część całego systemu bezzałogowego (UAS – \textit{Unmanned Aircraft System}). 

\paragraph{}
Który jest obecnie jednym z najszybciej rozwijających się segmentów nowoczesnych wojskowych systemów lotniczych. 
Wynika to przede wszystkim z niskich kosztów ich produkcji i eksploatacji w porównaniu z platformami załogowymi. 
Jednakże do ich największych zalet należy możliwość prowadzenia długotrwałych misji bez narażania życia załogi w strefach wysokiego ryzyka operacyjnego.

\paragraph{}
Wyżej wymienione systemy, składające się z bezzałogowego statku powietrznego (UAV), naziemnej stacji kontroli (GCS – \textit{Ground Control Station}) oraz naziemnego terminalu danych (GDT – \textit{Ground Data Terminal}).
Są projektowane do realizacji ściśle określonych zadań operacyjnych:
\begin{itemize}
    \item ISR \textit{Intelligence, Surveillance, and Reconnaissance} (wywiad, obserwacja i rozpoznanie)
    \item ISTAR \textit{Intelligence, Surveillance, Target Acquisition, Reconnaissance} - wywiad, obserwacja, idetyfikacjia celu, rozpoznanie
    \item CAS \textit{Close Air Support} - bliskie wsparcie powietrzn
    \item SEAD \textit{Suppression of Enemy Air Defenses} - przełamywanie obrony powietrznej wroga
    \item EW \textit{Electronic Warfare} - walka elektroniczna oraz zakłócanie
    \item C2 \textit{Command and Control} - dowodzenie i kontrola
    \item C4ISR \textit{Command, Control, Communications, Computers, Intelligence, Surveillance, and Reconnaissance} - dowodzenie, kontrola, łączność, komputery, wywiad, obserwacja i rozpoznanie
\end{itemize}

\paragraph{}
Prawidłowa identyfikacja wymagań i potrzeb klienta na wstępnym etapie jest kluczowa!
Ponieważ wprowadzanie zmian w profilu misji lub dodawanie nowych funkcjonalności
na późniejszych etapach projektu może wiąże się z koniecznością przebudowy całego projektu. 
Co drastycznie podnosi koszty oraz czas realizacji.
